\section*{अध्याय \ref{ch:g10:synthesis-yield} --- संश्लेषण: लब्धि और शुद्धता}

\begin{solution}{exo:g10:synthesis-yield:1}
$\eta = 4.2 / 5.0 = 0.84$, यानी \qty{84}{\percent} की \omterm{def:g10:synthesis-yield:yield}{लब्धि}।
\end{solution}

\begin{solution}{exo:g10:synthesis-yield:2}
\omterm{def:g7:chemical-reactions:reactant}{अभिकारकों} को तौलना; \omterm{def:g10:synthesis-yield:reflux}{पश्चवाह तापन}; \omterm{def:g10:synthesis-yield:synthesis}{अपरिष्कृत उत्पाद} का \omterm{def:g10:synthesis-yield:vacuum-filtration}{निर्वात
निस्यंदन}; \omterm{def:g10:synthesis-yield:recrystallisation}{पुनःक्रिस्टलन}; गलनांक मापना।
\end{solution}

\begin{solution}{exo:g10:synthesis-yield:3}
नीचे से डाला गया पानी ऊपर से निकलने से पहले पूरा आवरण भर देता है, इसलिए
संघनित्र की पूरी लंबाई ठंडी होती है। ऊपरी सिरा खुला छोड़ा जाता है ताकि उपकरण
बंद न हो: बंद उपकरण को गरम करने से दाब बनेगा और वह
फट सकता है।
\end{solution}

\begin{solution}{exo:g10:synthesis-yield:4}
उत्पाद गरम \omterm{def:g6:solutions-and-solubility:solute}{विलायक} में बहुत \omterm{def:g6:solutions-and-solubility:solute}{विलेय} और ठंडे \omterm{def:g6:solutions-and-solubility:solute}{विलायक} में लगभग अविलेय होना
चाहिए। अशुद्धियाँ, जो थोड़ी मात्रा में हैं, ठंडे \omterm{def:g6:solutions-and-solubility:solute}{विलायक} में
घुली रहती हैं और \omterm{def:g3:separating-mixtures:filter}{निस्यंद} में पहुँचती हैं।
\end{solution}

\begin{solution}{exo:g10:synthesis-yield:5}
नहीं: वह \qty{135}{\celsius} से नीचे और कई डिग्री के परास में पिघलता है, जो
अशुद्ध ठोस का संकेत है। उसका \omterm{def:g10:synthesis-yield:recrystallisation}{पुनःक्रिस्टलन} करना, उसे सुखाना,
और उसका गलनांक फिर से मापना चाहिए।
\end{solution}

\begin{solution}{exo:g10:synthesis-yield:6}
$n = 2.00 / 138.0 = \qty{0.0145}{mol}$ सैलिसिलिक अम्ल, जो \omterm{def:g10:reaction-progress-table:limiting}{सीमांत
अभिकारक} है; $m_{\max} = 0.0145 \times 180.0 = \qty{2.61}{g}$;
$\eta = 2.03 / 2.61 = 0.78$, यानी \qty{78}{\percent}।
\end{solution}

\begin{solution}{exo:g10:synthesis-yield:7}
तौले गए द्रव्यमान में पानी शामिल है, इसलिए वह ऐस्पिरिन के द्रव्यमान से अधिक है:
निकाली गई \omterm{def:g10:synthesis-yield:yield}{लब्धि} बहुत अधिक आती है, यहाँ असंभव रूप से
\qty{100}{\percent} से भी ऊपर। सूखे \omterm{def:g10:synthesis-yield:synthesis}{अपरिष्कृत उत्पाद} में अब भी अशुद्धियाँ हैं
(अभिक्रिया न किया सैलिसिलिक अम्ल, उपोत्पाद), जो द्रव्यमान बढ़ाती हैं: उसकी \omterm{def:g10:synthesis-yield:yield}{लब्धि} भी
शुद्ध ऐस्पिरिन की वास्तविक \omterm{def:g10:synthesis-yield:yield}{लब्धि} की तुलना में बहुत अधिक आती है।
\end{solution}

\begin{solution}{exo:g10:synthesis-yield:8}
यह कहीं तेज़ है, और क्रिस्टलों को कहीं अधिक सूखा छोड़ता है, क्योंकि
चूषण ज़्यादातर द्रव खींच लेता है; क्रिस्टलों को कीप पर ही धोना और चपटे
फ़िल्टर पेपर से खुरचकर निकालना भी आसान होता है।
\end{solution}

\begin{solution}{exo:g10:synthesis-yield:9}
उत्पाद $3.1/5.0 = 0.62$; ऐस्पिरिन $3.1/5.0 = 0.62$; सैलिसिलिक अम्ल
$2.2/5.0 = 0.44$। उत्पाद ऐस्पिरिन की तरह चलता है और सैलिसिलिक अम्ल का कोई धब्बा
नहीं दिखाता: वह ऐस्पिरिन है, और कोई सैलिसिलिक अम्ल नहीं पाया गया।
\end{solution}

\begin{solution}{exo:g10:synthesis-yield:10}
\omterm{def:g6:solutions-and-solubility:solute}{विलायक} A: गरम में बहुत \omterm{def:g6:solutions-and-solubility:solute}{विलेय}, ठंडे में लगभग अविलेय; \omterm{def:g6:solutions-and-solubility:solute}{विलायक} B ठंडा होने पर भी
ज़्यादातर उत्पाद को घुला रखता है। A के साथ \qty{3.0}{g} के लिए लगभग
$3.0 / 150 = \qty{0.020}{L}$, यानी \qty{20}{mL} गरम \omterm{def:g6:solutions-and-solubility:solute}{विलायक} चाहिए; ठंडा
होने पर अधिक से अधिक $2 \times 0.020 = \qty{0.04}{g}$ घुला रहता है।
\end{solution}

\begin{solution}{exo:g10:synthesis-yield:11}
उबलते द्रव के ऊपर जगह चाहिए: झाग और छींटे संघनित्र तक नहीं पहुँचने
चाहिए। क्वथन-पत्थर वाष्प को बुलबुले बनाने के लिए छोटी-छोटी जगहें देते हैं,
ताकि द्रव अचानक उफानों के बजाय एकसमान रूप से उबले।
\end{solution}

\begin{solution}{exo:g10:synthesis-yield:12}
$0.80 \times 0.70 = 0.56$: \qty{56}{\percent}। \qty{90}{\percent} के तीन
चरण: $0.90^3 = 0.73$, \qty{73}{\percent}। हर चरण में उत्पाद (और समय और पैसा)
खोता है, और हानियाँ गुणा होती जाती हैं: चरण जितने कम, अंत में उतना अधिक
उत्पाद बचता है।
\end{solution}

\begin{solution}{exo:g10:synthesis-yield:13}
अधिक से अधिक $3.3 \times 0.050 = \qty{0.17}{g}$ ऐस्पिरिन। बर्फ़-ठंडा पानी
ऐस्पिरिन को और भी कम घोलता है, क्योंकि उसकी \omterm{def:g6:solutions-and-solubility:solubility}{विलेयता} ताप के साथ
घटती है।
\end{solution}

\begin{solution}{exo:g10:synthesis-yield:14}
\begin{enumerate}
  \item $n_0(\text{सैलिसिलिक अम्ल}) = 2.76 / 138.0 = \qty{0.0200}{mol}$,
        ऐनहाइड्राइड \qty{0.0500}{mol}: सैलिसिलिक अम्ल सीमांत है,
        $x_{\max} = \qty{0.0200}{mol}$, और $0.0500 - 0.0200 =
        \qty{0.0300}{mol}$ ऐनहाइड्राइड बचता है।
  \item \omterm{def:g10:synthesis-yield:synthesis}{संश्लेषण} \qty{0.0200}{mol} एथेनोइक अम्ल देता है, आधिक्य का
        विनाश $2 \times 0.0300 = \qty{0.0600}{mol}$:
        कुल \qty{0.0800}{mol}।
\end{enumerate}
\end{solution}

\begin{solution}{exo:g10:synthesis-yield:15}
$n(\text{ऐस्पिरिन}) = \num{1.00e6} / 180.0 = \qty{5.56e3}{mol}$।
\qty{85}{\percent} की \omterm{def:g10:synthesis-yield:yield}{लब्धि} के साथ आवश्यक सैलिसिलिक अम्ल:
$\num{5.56e3} / 0.85 = \qty{6.54e3}{mol}$, यानी
$\num{6.54e3} \times 138.0 = \qty{9.0e5}{g}$, लगभग \qty{0.90}{t}।
\end{solution}

\begin{solution}{pb:g10:synthesis-yield:1}
\textbf{1.} C: $7 + 4 = 11$ और $9 + 2 = 11$; H: $6 + 6 = 12$ और
$8 + 4 = 12$; O: $3 + 3 = 6$ और $4 + 2 = 6$। संतुलित।

\textbf{2.} सैलिसिलिक अम्ल $7 \times 12.0 + 6 \times 1.0 + 3 \times 16.0
= \qty{138.0}{g/mol}$; ऐनहाइड्राइड $4 \times 12.0 + 6 \times 1.0 +
3 \times 16.0 = \qty{102.0}{g/mol}$; ऐस्पिरिन $9 \times 12.0 + 8 \times
1.0 + 4 \times 16.0 = \qty{180.0}{g/mol}$; एथेनोइक अम्ल $2 \times 12.0
+ 4 \times 1.0 + 2 \times 16.0 = \qty{60.0}{g/mol}$।

\textbf{3.} \omterm{def:g7:chemical-reactions:reactant}{अभिकारक}: सैलिसिलिक अम्ल और एथेनोइक ऐनहाइड्राइड; वांछित
उत्पाद: ऐस्पिरिन; उपोत्पाद: एथेनोइक अम्ल।

\textbf{4.} गरम करने से अभिक्रिया तेज़ होती है; पश्चवाह में वाष्प
संघनित होकर लौट आती है, इसलिए कोई \omterm{def:g7:chemical-reactions:reactant}{अभिकारक} या उत्पाद नहीं खोता।

\textbf{5.} $n = 3.0 / 138.0 = \qty{0.0217}{mol}$।

\textbf{6.} $m = 6.0 \times 1.08 = \qty{6.48}{g}$;
$n = 6.48 / 102.0 = \qty{0.0635}{mol}$।

\textbf{7.} सैलिसिलिक अम्ल: $0.0217 - x$; ऐनहाइड्राइड $0.0635 - x$;
ऐस्पिरिन और एथेनोइक अम्ल: $x$। सैलिसिलिक अम्ल पहले ख़त्म होता है:
$x_{\max} = \qty{0.0217}{mol}$, सैलिसिलिक अम्ल सीमांत है।

\textbf{8.} $m_{\max} = 0.0217 \times 180.0 = \qty{3.91}{g}$।

\textbf{9.} ताकि सारा सैलिसिलिक अम्ल ऐस्पिरिन में बदल जाए: जो बचता वह ऐस्पिरिन
में मिला रहता और उसे हटाना कठिन होता, जबकि आधिक्य ऐनहाइड्राइड अंत में मिलाए गए पानी
से नष्ट हो जाता है, और उससे बना एथेनोइक अम्ल धुलकर
निकल जाता है।

\textbf{10.} उसमें पानी शामिल है: \qty{3.6}{g} एक ग़लत
$3.6 / 3.91 = \qty{92}{\percent}$ देगा।

\textbf{11.} $3.3 / 3.91 = 0.84$: \qty{84}{\percent} सूखा \omterm{def:g10:synthesis-yield:synthesis}{अपरिष्कृत
उत्पाद}, जो शुद्ध नहीं है।

\textbf{12.} अशुद्धियाँ हट जाती हैं, और ऐस्पिरिन का कुछ भाग ठंडे \omterm{def:g6:solutions-and-solubility:solute}{विलायक} और
धोने के पानी में घुला रह जाता है।

\textbf{13.} अधिक से अधिक $3.3 \times 0.030 = \qty{0.10}{g}$।

\textbf{14.} $\eta = 2.7 / 3.91 = 0.69$: \qty{69}{\percent}।

\textbf{15.} शुद्ध ऐस्पिरिन वाले मान पर तीक्ष्ण गलनांक: क्रिस्टल
ऐस्पिरिन हैं, और शुद्ध हैं।

\textbf{16.} \qty{135}{\celsius} से नीचे और चौड़े परास में पिघलना:
\omterm{def:g10:synthesis-yield:synthesis}{अपरिष्कृत उत्पाद} अशुद्ध था।

\textbf{17.} \omterm{def:g10:synthesis-yield:synthesis}{अपरिष्कृत उत्पाद} में ऐस्पिरिन और अभिक्रिया न किया सैलिसिलिक
अम्ल था; \omterm{def:g10:synthesis-yield:recrystallisation}{पुनःक्रिस्टलन} ने सैलिसिलिक अम्ल हटा दिया।

\textbf{18.} $R_f = 3.1 / 5.0 = 0.62$।

\textbf{19.} दोनों में लगभग बराबर: संभव \qty{3.91}{g} से
\qty{3.3}{g} \omterm{def:g10:synthesis-yield:synthesis}{अपरिष्कृत उत्पाद} तक, और \omterm{def:g10:synthesis-yield:recrystallisation}{पुनःक्रिस्टलन} में \qty{3.3}{g} से \qty{2.7}{g} तक,
हर एक में लगभग \qty{0.6}{g}। चूँकि \omterm{def:g10:synthesis-yield:synthesis}{अपरिष्कृत उत्पाद} शुद्ध ऐस्पिरिन नहीं था,
काम के पहले भाग में वास्तव में \qty{0.6}{g} से थोड़ी अधिक
ऐस्पिरिन खोई।

\textbf{20.} \omterm{def:g10:synthesis-yield:synthesis}{संश्लेषण} की \omterm{def:g10:synthesis-yield:yield}{लब्धि} लगभग \qty{69}{\percent} है।
\end{solution}
